% GENERATED FILE. Edit canonical lessons, then run npm run generate:books.
% canonical-source-hash: fnv1a64:c873ee0828d10d9b
% canonical-lessons: ML-C62-rope, ML-C62-thread, ML-C62-needle, ML-C62-broom, ML-C62-winnowing-tray

\chapter{Made by Hand}
\label{ch:ml-made-by-hand}
\hlchaptermodality{62}

\begin{chapteropening}
\textbf{By the end of this chapter:} \emph{I can name rope, thread, a needle, a broom and a winnowing tray in Malayalam.}

Complete ML-C62-winnowing-tray and say all five words in order.
\end{chapteropening}

\section[kayaṟ]{\ml{കയർ} (kayaṟ) — coir rope}

\label{lesson:ML-C62-rope}



\begin{quote}
Before the new run of words: what did \emph{śarkkara} mean, and what did \emph{muḷakŭ} mean?
\end{quote}

\subsection*{You'll want to know: \ml{കയർ}}
\textbf{\ml{കയർ}} (\emph{kayaṟ}) — \textquotedblleft{}coir rope\textquotedblright{}.

Native Dravidian, from a root about twisting, and it names both the fibre and the rope twisted out of it.

The fibre comes off the husk of the coconut you met on the shelf, soaked in backwater until it softens and then beaten out by hand. Kerala has made rope this way for as long as anyone has written about it.

English borrowed the word along with the trade: \emph{coir} is this word, carried out of Malabar on Portuguese and Dutch ships.

\begin{grammarlens}[title={What you've built}]
The first of five things made by hand.
\end{grammarlens}

\subsection*{Your turn}
Say these aloud:

\begin{itemize}
\raggedright
  \item \emph{kayaṟ}
  \item it once more, slowly
  \item \emph{tēṅṅa}, then \emph{kayaṟ}, and say which one the other is made from
\end{itemize}

\begin{tcolorbox}[breakable,colback=teal!4,colframe=teal!35!black,title={Before you move on}]
What does \ml{കയർ} mean? (\textquotedblleft{}coir rope\textquotedblright{}.) Say it once more.
\end{tcolorbox}
\section[nūlŭ]{\ml{നൂല്} (nūlŭ) — thread}

\label{lesson:ML-C62-thread}



\begin{quote}
Before the new one: what did \emph{kayaṟ} mean?
\end{quote}

\subsection*{You'll want to know: \ml{നൂല്}}
\textbf{\ml{നൂല്}} (\emph{nūlŭ}) — \textquotedblleft{}thread\textquotedblright{}.

Native Dravidian, matching Tamil \emph{nūl}, and it stands behind \emph{tuṇi}, the cloth-word: one is what the other is woven from.

The older language put the same word to a second job. A treatise, a work of scholarship, was also a \emph{nūl} — a thread of argument running from one end of a subject to the other, which is close to what English does with a train of thought.

It closes on the bare \emph{l} of \emph{cōṟu} and \emph{vērŭ}: the consonant with its vowel shaved off.

\begin{grammarlens}[title={What you've built}]
Two.
\end{grammarlens}

\subsection*{Your turn}
Say these aloud:

\begin{itemize}
\raggedright
  \item \emph{nūlŭ}
  \item it once more, slowly
  \item \emph{kayaṟ}, then \emph{nūlŭ}, and say which of the two is thicker
\end{itemize}

\begin{tcolorbox}[breakable,colback=teal!4,colframe=teal!35!black,title={Before you move on}]
What does \ml{നൂല്} mean? (\textquotedblleft{}thread\textquotedblright{}.) Say it once more.
\end{tcolorbox}
\section[sūci]{\ml{സൂചി} (sūci) — a needle}

\label{lesson:ML-C62-needle}



\begin{quote}
Before the new one: what did \emph{nūlŭ} mean?
\end{quote}

\subsection*{You'll want to know: \ml{സൂചി}}
\textbf{\ml{സൂചി}} (\emph{sūci}) — \textquotedblleft{}a needle\textquotedblright{}.

Sanskrit \emph{sūcī}, from a root about piercing and pointing, and the same root gives Sanskrit its word for an indication or a sign.

So the needle and the pointer come from one idea: the thing that goes in, and the thing that points out. A Malayalam newsreader saying that a report \emph{indicates} something is standing on this root.

With \emph{nūlŭ} through the eye of it, a \ml{സൂചി} is what turns cloth back into something a person can wear.

\begin{grammarlens}[title={What you've built}]
Three.
\end{grammarlens}

\subsection*{Your turn}
Say these aloud:

\begin{itemize}
\raggedright
  \item \emph{sūci}
  \item it once more, slowly
  \item \emph{nūlŭ}, then \emph{sūci}, and say which one goes through the other
\end{itemize}

\begin{tcolorbox}[breakable,colback=teal!4,colframe=teal!35!black,title={Before you move on}]
What does \ml{സൂചി} mean? (\textquotedblleft{}a needle\textquotedblright{}.) Say it once more.
\end{tcolorbox}
\section[cūlŭ]{\ml{ചൂല്} (cūlŭ) — a broom}

\label{lesson:ML-C62-broom}



\begin{quote}
Before the new one: what did \emph{sūci} mean?
\end{quote}

\subsection*{You'll want to know: \ml{ചൂല്}}
\textbf{\ml{ചൂല്}} (\emph{cūlŭ}) — \textquotedblleft{}a broom\textquotedblright{}.

Native Dravidian, and a Kerala \ml{ചൂല്} is a bundle of coconut-leaf ribs bound tight at one end — the coconut again, a third time on a third object.

Sweeping the yard at first light is the opening work of a Kerala day, done before the kitchen is lit, and it leaves the ground marked in long strokes that stay until somebody walks on them.

There is a second broom for indoors, softer, made of grass rather than leaf-rib, and the same word covers it.

\begin{grammarlens}[title={What you've built}]
Four.
\end{grammarlens}

\subsection*{Your turn}
Say these aloud:

\begin{itemize}
\raggedright
  \item \emph{cūlŭ}
  \item it once more, slowly
  \item \emph{sūci}, then \emph{cūlŭ}, and say which of the two you hold in a fist
\end{itemize}

\begin{tcolorbox}[breakable,colback=teal!4,colframe=teal!35!black,title={Before you move on}]
What does \ml{ചൂല്} mean? (\textquotedblleft{}a broom\textquotedblright{}.) Say it once more.
\end{tcolorbox}
\section[muṟaṁ]{\ml{മുറം} (muṟaṁ) — a winnowing tray}

\label{lesson:ML-C62-winnowing-tray}



\begin{quote}
Before the new one: what did \emph{cūlŭ} mean?
\end{quote}

\subsection*{You'll want to know: \ml{മുറം}}
\textbf{\ml{മുറം}} (\emph{muṟaṁ}) — \textquotedblleft{}a winnowing tray\textquotedblright{}.

Native Dravidian: a flat plaited tray with a lip on three sides, woven from bamboo strip like the mat and the basket before it.

Its work needs the wind. Grain is tossed up off the \ml{മുറം} and \emph{kāṟṟŭ} carries the husk sideways while the heavy part drops back — a machine with no moving parts, run on weather.

\emph{Ari}, the raw rice on the kitchen shelf, has been through one of these on its way there.

\begin{grammarlens}[title={What you've built}]
Five, and every one of them came off somebody's hands: rope, thread, a needle, a broom and a winnowing tray.
\end{grammarlens}

\subsection*{Your turn}
Say these aloud:

\begin{itemize}
\raggedright
  \item \emph{muṟaṁ}
  \item it once more, slowly
  \item all five in order — \emph{kayaṟ}, \emph{nūlŭ}, \emph{sūci}, \emph{cūlŭ}, \emph{muṟaṁ}
\end{itemize}

\begin{tcolorbox}[breakable,colback=teal!4,colframe=teal!35!black,title={Before you move on}]
What does \ml{മുറം} mean? (\textquotedblleft{}a winnowing tray\textquotedblright{}.) Say it once more.
\end{tcolorbox}
